\documentclass{beamer}
\usepackage{comment}
\title{MATH 1190 \& 1210 Meeting for Week 3}
\author{Josh}
\date{26 January 2026}

\newcommand{\transition}[1]{\frame{\begin{center}\fbox{#1}\end{center}}}

\begin{document}

\frame{\titlepage}

\section{Logistics}

\transition{Logistics}

\frame[t]{{Logistics}

\begin{itemize}
\item please update your office hours \& LA name on Canvas {\bf ASAP}
    \begin{itemize}
    \item 3 hours/week, at least 1 hour/week in the MCLC
    \item let me know if you want to use Gilmer 490 (Academic Commons) - we can reach out to the scheduling folks to inquire
    \item please minimize overlap with other instructors' office hours to maximize availability for students
    \item update LA's name; if you're not sure what your LA's name is, then your homework is to learn your LA's name :) 
    \end{itemize}
\item Exam 1
    \begin{itemize}
    \item Exam 1 is Thursday Feb 5 (Next week).
    \item Please look over the exam and see how the problems feel.
    \item I will make homework for lesson 5 due on Tuesday February 3rd so that students have time to ask instructors questions.
    \item Please fill out the proctoring sign up sheet if you have not already.
    \item I will send discord invites and grading assignments by the end of the week. 
    \item Please create a rubric and write up a solution to your problems on overleaf. Then share the rubric on discord so we may discuss.
    \end{itemize}
\end{itemize}
}

\frame[t]{{Logistics}

\begin{itemize}
\item deadlines
    \begin{itemize}
    \item if you find you are behind or ahead of schedule, then feel free to shift deadlines of pre-class assignments for your section(s) accordingly
    \item select the assignment, click "Edit Assignment Settings", then click "Manage Due Dates and Assign To", then find your section and change the time as needed
    \item pre-class assignments should always be due at the start of the class during which the corresponding activity is to be introduced 
    \end{itemize}
\item attendance
    \begin{itemize}
    \item not a big deal if you haven't been recording it since there's lots of movement into/out of classes right now
    \item take attendance everyday starting this week, sooner is better though to normalize attendance checking if you want to add it as a grade we will place it under the ``activities" category
    \end{itemize}
\end{itemize}

}

\frame[t]{{Logistics}

\begin{itemize}
\item Room Assignment
    \begin{itemize}
    \item will occur during week 3 after last day to drop (1/27)
   % \item talking about the curve $\to$ fewer student grumbles
    \item please be prepared to answer student questions about letter grade assignment 
    \end{itemize}
\item in-class activities
    \begin{itemize}
    \item customizable; up to 10 points per
    \item just create assignments in Canvas for your section as needed
    \item no more than 10 points per activity toward their activities grade; for example, we use 10-points per activity in MATH 1190 discussions, activities tend to take most if not all of the 50-minute class period
    \end{itemize}
\end{itemize}
}

\frame{{Logistics}

\begin{itemize}
\item student questions
    \begin{itemize}
    \item each instructor should be the primary point of contact for their own students when it comes to student questions
    %\item please don't forward student emails directly to DJ, or ask students to contact him directly (unless they are one of our students)
    \item if a student asks a question you're unsure of, then post it in discord -- chances are someone else has/had the same question, and perhaps can help!
    %\item always feel free to ask how to answer a student's questions if you need any guidance with this - I am happy to help! 
    \end{itemize}
\item Active learning activities
    \begin{itemize}
    \item I will send a sign up sheet for you to discuss your active learning exercises later this week
    \item intention is to provide support in your growth as an educator
    \end{itemize}
\end{itemize}
}

\section{Content}

\transition{Content}

\frame[t]{{Content}

{\bf Limits \& Infinity}
\begin{itemize}
\item dense lesson, will probably take about 100 minutes in 1190 
\item pre-class assignment
    \begin{itemize}
    \item "when $x$ is very large, what happens to $f(x)$?
    \item constant over something growing without bound $\implies$ zero limit
    \item some functions grow faster than others: $e^x$ grows faster than $x^2$, which grows faster than $\ln(x)$
    \item intro to forms $\frac{c}{\infty}$ \& $\frac{\infty}{\infty}$  
    \end{itemize}
\end{itemize}
}
\begin{comment}
\frame[t]{{Content}
\begin{itemize}
\item in-class lesson
    \begin{itemize}
    \item Warm-up Exercise 1 picks up where pre-class left off
    \end{itemize}
\end{itemize}

\includegraphics[width=\textwidth]{images/lesson4warmupexercise.png}

}

\frame[t]{{Content}
\begin{itemize}
\item mini-lecture, pg. 1
    \begin{itemize}
    \item students need to know which kinds of functions have infinite limits at infinity
    \item might ask students to conjecture which kinds of functions have infinite limits at infinity, as they'll have thought about this in the pre-class / can look at the given graphs
    \end{itemize}
\end{itemize}

\includegraphics[width=\textwidth]{images/lesson4minilecture1.png}

}

\frame[t]{{Content}

{\bf Limits \& Infinity}
\begin{itemize}
    \item Poll / Example 1 / Exercise A: hypothesize in the poll, then test out hypothesis in example/exercise; constant over something growing without bound $\implies$ zero limit
    \begin{itemize}
    \item Can (should?) convert some of Example 1 into an exercise, as students will have done some of these problems in the pre-class.
    \end{itemize}
\end{itemize}

\includegraphics[width=\textwidth]{images/lesson4example1.png}

}

\frame[t]{{Content}

\begin{itemize}
    \item Poll / Example 2 / Exercise B: hypothesize in the poll, then test out hypothesis in example/exercise; something growing large over constant $\implies$ infinite limit, need to check signs
    \begin{itemize}
    \item this form is not discussed in the pre-class assignment $\implies$ maybe less safe to convert some of the example to an exercise
    \end{itemize}
\end{itemize}

\includegraphics[width=\textwidth]{images/lesson4example2.png}

}

\frame[t]{{Content}

\begin{itemize}
    \item Poll / Example 3 / Exercise C: hypothesize in the poll, then test out hypothesis in example/exercise; ``pecking order" / ``dominance hierarchy"; something growing large over something growing large $\implies$ compare terms that dictate growth for large input values, refer back to previous cases
    \begin{itemize}
        \item students will have seen that if the faster-growing term is in the denominator, then the limit has a value of 0, but will not be familiar with other cases
    \end{itemize}
\end{itemize}

\

\includegraphics[width=\textwidth]{images/lesson4example3.png}

}

\frame[t]{{Content}

\begin{itemize}
    \item Poll / Example 4 / Exercise D: hypothesize in the poll, then test out hypothesis in example/exercise; constant over something becoming closer and closer to 0 $\implies$ DNE, possibly can say $=\infty$ or $=-\infty$
        \begin{itemize}
        \item these are challenging for students, as they tend to think that the sign on the infinity is inherited from the superscript in the limit
        \end{itemize}
\end{itemize}

\

\,\hfill $\displaystyle\lim_{x\to1^{-}} \frac{\sqrt{x}+1}{x-1}$ \hfill $\displaystyle\lim_{x\to3} \frac{x^2-2x}{3-x}$ \hfill $\displaystyle\lim_{x\to 0}\frac{1+x}{x^2}$ \hfill\,

}
\end{comment}

\frame{{Content}

{\bf Derivatives \& Rates of Change}
\begin{itemize}
\item will probably take about 60 minutes in 1190
\item pre-class assignment (H5P + written)
    \begin{itemize}
    \item finding units of a rate of change
    \item AROC $\approx$ IROC, and approx improves as $\Delta x\to0$
    \item Warm-up Exercise A: infer IROC/tangent line slope is defined by a limit of AROC/secant line slope 
    \end{itemize}
\item in-class lesson
    \begin{itemize}
    \item Mini-lecture: derivative is defined by a (1) limit of (2) AROC/secant line slopes, and therefore (1) may sometimes not exist but (2) represents IROC/tangent line slopes when it does exist 
    \item Example 1/Exercise B: derivative as slope 
    \item Exercise C: derivative as slope; height of derivative's graph is slope of function's graph; higher CD version of Exercise B 
    \item Exercise D: difference between derivative and difference quotient; derivative as rate of change; units of measurement; marginal cost
    \end{itemize}
\end{itemize}
}



%------------------------------------------------
\section{Pedagogy}

\transition{Pedagogy}
%\frame[t]{Pedagogy}

%------------------------------------------------
\begin{frame}{Why Rubrics Matter in Coordinated Courses}
\begin{itemize}
    \item Rubrics communicate clear expectations to students
    \item Promote fairness and consistency across sections and instructors
    \item Improve grading efficiency, especially in large-enrollment calculus courses
    \item Provide meaningful feedback beyond numerical scores
\end{itemize}

\vspace{0.3cm}
\textbf{Key idea:} Rubrics support both instructional alignment and equity.
\end{frame}

%------------------------------------------------
\begin{frame}{What Is a Rubric?}
A rubric is an evaluation tool that:
\begin{itemize}
    \item Defines criteria for an assignment or learning outcome
    \item Describes levels of achievement in clear, measurable terms
    \item Can be used for exams, quizzes, homework, projects, and presentations
\end{itemize}

\vspace{0.3cm}
In calculus, rubrics are especially useful for:
\begin{itemize}
    \item Multi-step problems
    \item Written explanations and proofs
    \item Modeling, applications, and group work
\end{itemize}
\end{frame}

%------------------------------------------------
\begin{frame}{Step 1: Clarify Purpose and Context}
Before creating a rubric, ask:
\begin{itemize}
    \item What learning outcomes are being assessed?
    \item Is this assessment formative (practice) or summative (final)?
    \item Will students revise their work based on feedback?
\end{itemize}

\vspace{0.3cm}
For coordinated calculus exams:
\begin{itemize}
    \item Are we assessing conceptual understanding, procedural fluency, or both?
    \item What does an ``excellent'' solution look like?
\end{itemize}
\end{frame}

%------------------------------------------------
\begin{comment}
\begin{frame}{Step 2: Choosing a Rubric Type}
Three common rubric types:
\begin{itemize}
    \item \textbf{Holistic} – one overall judgment
    \item \textbf{Analytic} – multiple criteria scored separately
    \item \textbf{Single-point} – focus on proficiency with narrative feedback
\end{itemize}

\vspace{0.3cm}
Choice depends on:
\begin{itemize}
    \item Assessment goals
    \item Time constraints
    \item Desired level of feedback
\end{itemize}
\end{frame}

%------------------------------------------------
\begin{frame}{Holistic Rubrics}
\textbf{Description:}
\begin{itemize}
    \item Assigns one overall score based on broad descriptions
    \item Emphasizes overall quality of work
\end{itemize}

\textbf{Advantages:}
\begin{itemize}
    \item Faster grading
    \item Useful for final exams or cumulative assessments
\end{itemize}

\textbf{Limitations:}
\begin{itemize}
    \item Limited diagnostic feedback
    \item Less effective for uneven student work
\end{itemize}
\end{frame}
\end{comment}
%------------------------------------------------
\begin{frame}{Analytic Rubrics}
\textbf{Description:}
\begin{itemize}
    \item Breaks work into criteria (e.g., setup, computation, interpretation)
    \item Each criterion is scored separately
\end{itemize}

\textbf{Advantages:}
\begin{itemize}
    \item Transparent partial credit
    \item Supports consistent grading across multiple graders
    \item Allows weighting (e.g., concepts vs. algebra)
\end{itemize}

\textbf{Best fit for calculus:}
\begin{itemize}
    \item Exams, quizzes, written solutions
\end{itemize}
\end{frame}

%------------------------------------------------
\begin{frame}{Single-Point Rubrics}
\textbf{Description:}
\begin{itemize}
    \item Defines only the ``proficient'' level
    \item Instructors add comments for strengths and weaknesses
\end{itemize}

\textbf{Advantages:}
\begin{itemize}
    \item Encourages growth-oriented feedback
    \item Reduces student focus on point-chasing
\end{itemize}

%\textbf{Best uses in calculus:}
%\begin{itemize}
%    \item Projects, modeling tasks, written explanations
%\end{itemize}
\end{frame}

%------------------------------------------------
\begin{frame}{Step 2: Starting Points for Rubric Design}
Options include:
\begin{itemize}
    \item Adapting existing or previous 1190/1210 rubrics
    \item Using AI tools to generate a draft (with careful revision)
    \item Building from scratch based on learning outcomes
\end{itemize}

\vspace{0.3cm}
\textbf{Recommendation:}
Maintain a shared rubric repository for coordinated courses. (Discord/WebAssign)
\end{frame}

%------------------------------------------------
\begin{frame}{Step 3: Defining Grading Criteria}
Effective criteria:
\begin{itemize}
    \item Align directly with learning objectives
    \item Focus on observable features of student work
    \item Use clear, student-friendly language
\end{itemize}

\vspace{0.3cm}
Examples for 1190-1210:
\begin{itemize}
    \item Correct interpretation of the problem
    \item Appropriate mathematical setup
 %   \item Accuracy of computations
    \item Quality of explanation or justification
\end{itemize}
\end{frame}

%------------------------------------------------
\begin{frame}{Step 4: Describing Performance Levels}
Best practices:
\begin{itemize}
    \item Leave plenty of comments
    \item Use distinct, clearly differentiated levels
    \item Avoid vague terms like ``good'' or ``poor''
    \item Emphasize quality and reasoning, not just errors
\end{itemize}

\vspace{0.3cm}
In 1190-1210 grading:
\begin{itemize}
    \item Describe conceptual errors vs. minor algebra slips
    \item Focus on impact on understanding and interpretation
\end{itemize}
\end{frame}

%------------------------------------------------
\begin{frame}{Step 5: Piloting and Calibration}
Before full implementation:
\begin{itemize}
    \item Apply the rubric to sample student work
    \item Compare scores across instructors or GTAs
    \item Revise unclear criteria or descriptions
\end{itemize}

\vspace{0.3cm}
\textbf{For coordinated courses:}
Hold grading norming sessions before major exams. 
\end{frame}

%------------------------------------------------
\begin{frame}{Rubrics and Test Planning}
Rubrics support exam design by:
\begin{itemize}
    \item Clarifying how partial credit will be awarded
    \item Ensuring alignment between problems and learning goals
    \item Helping balance conceptual understanding and computation
\end{itemize}

\vspace{0.3cm}
Well-designed rubrics lead to:
\begin{itemize}
    \item Fewer grading disputes
    \item More actionable feedback
\end{itemize}
\end{frame}

%------------------------------------------------
\begin{frame}{Best Practices for Coordinated Calculus Courses}
\begin{itemize}
    \item Keep rubrics concise and readable
 %   \item Share rubrics with students before assessments
    \item Use consistent language across sections
    \item Check alignment with gradebook and exam criteria
    \item Check alignment with practice exams
    \item Ensure accessibility (screen-reader friendly formats)
\end{itemize}
\end{frame}

%------------------------------------------------
\begin{frame}{Discussion Questions}
\begin{enumerate}
    \item Where do graders most often disagree when scoring calculus exams?
%    \item Which calculus assessments would benefit most from analytic rubrics?
    \item How can rubrics help balance procedural fluency and conceptual understanding?
    \item What level of detail is appropriate for partial credit in multi-step problems?
    \item How can rubric calibration be better integrated into our discussions?
\end{enumerate}
\end{frame}

%------------------------------------------------
\begin{frame}{Source and Acknowledgment}
Content adapted from:

\vspace{0.3cm}
\small
NC State University DELTA Teaching Resources\\
\textit{Rubric Best Practices, Examples, and Templates}

\vspace{0.3cm}
\url{https://teaching-resources.delta.ncsu.edu/rubric_best-practices-examples-templates/}

\vspace{0.3cm}
Licensed under Creative Commons Attribution–NonCommercial 4.0.
\end{frame}

\end{document}

\end{document}


%%%%%%%%%%%%%%%%%%%%%%%%%%%%%%%%%%%

\section{Pedagogy}

\transition{Pedagogy}

\frame{{Pedagogy}

How can we view difficulty of a task from a students' perspective?\\[0.1in]

\textbf{\underline{Cognitive Demand}} 

({\footnotesize adapted from {\color{blue}\href{https://journals.sagepub.com/doi/abs/10.3102/00028312033002455}{Stein et al}}, 1996}) 

\vfill

$\text{low-level demand }\begin{cases} \text{(1) Memorization tasks}\\
\text{(2) Procedural tasks}\end{cases}$ 

\vfill

$\text{high-level demand }\begin{cases} \text{(3) Connections tasks}\\
\text{(4) Mathematics tasks}\end{cases}$ 

\vfill

$\text{no demand }\begin{cases} \text{(0) student disengaged}\end{cases}$ 

\vfill

A student can experience very different CD levels while engaging in the same task depending on how the task is scaffolded.

}

\frame{{Pedagogy}

Assuming the examples are worked first, what would the CD level be for a typical 1210 student on the exercises?

\begin{center}
\includegraphics[width=1.1\textwidth]{images/m1210lesson5exerciseB.PNG}
\end{center}

Recall: (1) memorization; (2) procedures w/o connections; (3) procedures w/ connections; (4) mathematics
}

\frame{{Pedagogy}

What would the CD level be for a typical 1210 student on this task, given the previous exercise had already been worked? {\bf Follow-up:} how can we scaffold this task beforehand (e.g. in exercise B) to help students be more successful in this task?

\begin{center}
\includegraphics[width=0.99\textwidth]{images/m1210lesson5exerciseC.PNG}
\end{center}

Recall: (1) memorization; (2) procedures w/o connections; (3) procedures w/ connections; (4) mathematics
}